\section{No charged instability above \texorpdfstring{$\alpha_\star = 2/3$}{alpha* = 2/3}}
\label{sec:4}

\subsection{At zero temperature the brane interior is \texorpdfstring{$\mathrm{AdS}_3 \times \mathbb{R}^2$}{AdS3 x R2}}
\label{sec:4.1}

Since the field is measured in units of $T^2$, cooling at fixed physical
field sends $B u_H^2$, and with it $\beta = \alpha (B u_H^2)^2$, to
infinity: the zero-temperature limit at fixed $\alpha$ is
$\beta \to \infty$. In this limit the interior of the brane develops a long
throat~\cite{DHoker:2009mmn}. It approaches the product geometry
$\mathrm{AdS}_3 \times \mathbb{R}^2$, in which the flux plane $(x,y)$ has
frozen at a finite proper size and the direction $z$ along the field has
joined $t$ and $r$ to form a three-dimensional anti-de Sitter space. In our
conventions $U = (r - r_0)^2/\ell_3^2$, $e^{2Z} = c_Z (r - r_0)^2$,
$e^{2V} = v$ solves the full Einstein--Yang--Mills equations if and only
if
\begin{equation}
\label{eq:throat}
\ell_3^2 = \tfrac13, \qquad v = B\sqrt{\alpha/6}~.
\end{equation}
The constants $c_Z$ and $r_0$ can be removed by rescaling $z$ and shifting
$r$.

Neither $v$ nor $B$ is meaningful on its own, since rescaling the flux
plane, $(x,y) \to k\,(x,y)$, sends both $e^{2V}$ and $B$ to themselves
divided by $k^2$. The scale-invariant form of~\eqref{eq:throat} is
\begin{equation}
\alpha B^2 e^{-4V} = 6, \qquad B\,e^{-2V} = \sqrt{6/\alpha}~.
\end{equation}
The first is the condition for the throat to exist. The second is the
magnetic field measured in the proper units of the frozen flux plane, and
it controls the condensate. At finite $\beta$ the throat joins the
$\mathrm{AdS}_5$ region at the crossover radius $r_c = (\beta/6)^{1/4}$,
where the outer growth $e^{2V} \simeq r^2$ meets the frozen value
$\sqrt{\beta/6}$ (figure~\ref{fig:throat}).

\begin{figure}[t]
\centering
\includegraphics{throat_geometry}
\caption{The magnetic brane forms its throat. (a)~The proper size $e^{2V}$
of the flux plane against $\ln r$ at fixed temperature, for
$\beta = 0, 1, 45, 10^4, 10^8$ (light to dark; $\beta = 0$ dashed). Outside
the crossover radius $r_c$ (circles) it grows as $r^2$, as in
$\mathrm{AdS}_5$; inside it the flux plane freezes at
$e^{2V} \to v = \sqrt{\beta/6}$. (b)~The invariant $\beta e^{-4V}$, which
is $0$ at the boundary and $6$ on the $\mathrm{AdS}_3 \times \mathbb{R}^2$
fixed point. Its horizon value is $3.58$, $5.613$ and $5.9877$ at
$\beta = 45$, $10^4$ and $10^8$, and the plateau near $6$ lengthens as
$\ln r_c = \tfrac14\ln(\beta/6)$.}
\label{fig:throat}
\end{figure}

At small temperature the throat is a BTZ black hole~\cite{Banados:1992wn}
times the flux plane:
$U = ((r-r_0)^2 - h^2)/\ell_3^2$, with horizon at $r - r_0 = h$, solves the
same equations with the same $\ell_3$ and $v$
(appendix~\ref{app:numerics-brane}). This is the throat on the
finite-temperature onset curve, and it governs the approach to
$\alpha_\star$ in section~\ref{sec:4.6}.

\subsection{The field drops out and the onset becomes a Breitenlohner--Freedman bound}
\label{sec:4.2}

On the throat metric the Sturm--Liouville equation of section~\ref{sec:3}
becomes an Euler equation in $\rho = r - r_0$, $(\rho^3 w')' + b\,\rho\,w = 0$, and $w = \rho^{-\Delta}$
gives the indicial equation
\begin{equation}
\Delta(\Delta - 2) = -b, \qquad b \equiv B\,\ell_3^2\,e^{-2V} = \sqrt{\frac{2}{3\alpha}}~.
\end{equation}
The field has cancelled because the frozen flux plane grows in
proportion to $B$, so the field in throat units, $b$, is fixed by the
coupling alone. In the language of the $\mathrm{AdS}_3$ factor, the lowest
Landau level of the $W$ boson is a scalar of mass $m^2\ell_3^2 = -b$, and
the Breitenlohner--Freedman
bound~\cite{Breitenlohner:1982bm,Breitenlohner:1982jf}
$m^2\ell_3^2 \ge -1$ is violated exactly
when $b > 1$. Violating it is sufficient for an instability at zero
temperature, and we show below that it is also necessary, so that
\begin{equation}
\text{the instability survives to } T = 0 \quad\Longleftrightarrow\quad \alpha < \alpha_\star = \tfrac23~.
\end{equation}
We call $\alpha_\star$ the critical coupling. As $\alpha \to 0$, $b \to \infty$
and we recover the probe tachyon of~\cite{Ammon:2011je}. Increasing
$\alpha$ lowers the throat field $b$, which squeezes the tachyon until, at
$\alpha_\star$, the two indicial roots meet at $\Delta = 1$.

That the bound is also necessary needs an argument, because
violating the bound of a throat is sufficient for an instability but not in
general necessary~\cite{Faulkner:2010gj,Horowitz:2010gk}: a bound state can
sit in the crossover between the throat and the boundary. For charged
scalars the two criteria were found numerically to
agree~\cite{Denef:2009tp}, and the bound is used as the criterion
in~\cite{Gubser:2009cg,Horowitz:2009ij}, but this is a
property of each model, and in $d = 3$ they disagree
(section~\ref{sec:4.3}).

Here the zero-temperature brane is a single geometry, and the coupling
enters the zero-mode problem only through $b = \sqrt{\alpha_\star/\alpha}$,
which plays the role of an eigenvalue. In the throat the mode behaves as
$(r - r_0)^{-1 \pm \sqrt{1 - b}}$, so the operator has a continuum for
$b > 1$, and a zero mode at a coupling above $\alpha_\star$ would be a
bound state below the threshold $b = 1$. Such a bound state and the
source-free solution at threshold have vanishing Wronskian at both ends,
at the boundary because both are source-free and in the throat because the
bound state decays, so by Sturm's comparison
theorem~\cite[theorem~9.39]{Teschl:2009mmqm} the threshold solution would
have a zero. At threshold the solution in the throat has the form
$(r - r_0)\,w = c_0 + c_1\ln\big((r - r_0)/\rho_\star\big)$, with
$\rho_\star$ the top of the throat (section~\ref{sec:4.6}), and a zero
inside the throat would show up as $c_1/c_0 > 0$. We find
no zero and $c_1/c_0 = -0.035$: the would-be zero lies far
outside the throat, at $\ln\big((r - r_0)/\rho_\star\big) = +28.6$. There is therefore
no bound state, and for every $\alpha > \alpha_\star$ the D'Hoker--Kraus
brane is stable against this mode at every field. A direct check agrees:
integrated from the throat to the boundary at couplings between
$\alpha_\star$ and $4\alpha_\star$, the solution that decays in the throat
always arrives with a source of the same sign
(appendix~\ref{app:numerics-onset}).

At finite temperature the statement is numerical. On the onset curve
$\alpha/\alpha_\star = (1 - \varepsilon/6)/b_h^2$, with
$\varepsilon = 6 - \beta e^{-4V}$ and $b_h = B_c\ell_3^2 e^{-2V}$ at the
horizon. Since $\varepsilon > 0$ at finite temperature, $b_h \ge 1$
implies $\alpha < \alpha_\star$. Shooting from the horizon up to
$\beta \approx 10^{66}$, where $\alpha = 0.6626$, the onset coupling rises
strictly and stays below $\alpha_\star$, with $b_h \ge 1.003$
(appendix~\ref{app:numerics-onset}); beyond that the matching analysis of
section~\ref{sec:4.6} takes over.

\subsection{In \texorpdfstring{$d+1$}{d+1} dimensions \texorpdfstring{$\alpha_\star = 16/(d(d-1)(d-2))$}{alpha* = 16/(d(d-1)(d-2))}}
\label{sec:4.3}

The argument is not specific to four boundary dimensions. The reduction to
the polarised lowest Landau level goes through as before, with
$P = U e^{(d-3)Z}$ and $Q = B e^{(d-3)Z - 2V}$; the $d - 3$ directions along
the field only contribute a volume factor. The throat is now
$\mathrm{AdS}_{d-1} \times \mathbb{R}^2$, one of the magnetic vacua that
exist in every dimension~\cite{Lu:2013eoa}, with
\begin{equation}
\ell_{d-1}^2 = \frac{(d-2)^2}{d(d-1)}, \qquad \frac{B\ell_{d-1}^2}{v} = \frac{(d-2)^{3/2}}{\sqrt{\alpha\,d(d-1)}}~,
\end{equation}
and the indicial equation becomes $\Delta(\Delta - (d-2)) = -B\ell_{d-1}^2/v$,
with $v = e^{2V}$ again the frozen size of the flux plane. Comparing with
the Breitenlohner--Freedman value $(d-2)^2/4$ for $\mathrm{AdS}_{d-1}$
gives
\begin{equation}
\alpha_\star(d) = \frac{16}{d(d-1)(d-2)}~,
\end{equation}
which is $8/3$, $2/3$, $4/15$ and $2/15$ for $d = 3, 4, 5, 6$.

For $d = 4, 5, 6$ this bound is the critical coupling. We constructed the
zero-temperature brane in each dimension, flowing from
$\mathrm{AdS}_{d+1}$ into the throat along its single irrelevant
deformation, and the threshold mode on it has no zero, with
$c_1/c_0 < 0$ (table~\ref{tab:dims}). At finite temperature the
onset coupling for $d = 5, 6$, as for $d = 4$, rises monotonically, stays
below $\alpha_\star(d)$ and approaches it as $T \to 0$
(appendix~\ref{app:numerics-onset}).

In $d = 3$ the bound is not the critical coupling. The throat is
$\mathrm{AdS}_2 \times \mathbb{R}^2$, and on the extremal magnetic
Reissner--Nordstr\"om--$\mathrm{AdS}_4$ brane the threshold mode has one
zero deep in the throat, so a bound state sits in the crossover below the
$\mathrm{AdS}_2$ bound. It disappears only at
\begin{equation}
\alpha_c(3) = 3.0370~,
\end{equation}
which we find by three methods and confirm with an independent
construction (appendix~\ref{app:numerics-onset}). At finite temperature the
onset coupling rises through $8/3$ and tends to $\alpha_c(3)$ as $T \to 0$.
The bound of~\cite{Wong:2013rda} on the magnetic $\mathrm{AdS}_4$ brane,
$\alpha_\star(3) = 8/3$ in our variables, is therefore sufficient for the
instability but, contrary to the statement in~\cite{Wong:2013rda} that the
$W$ bosons condense only for $\alpha < 8/3$, not necessary; there the
zero-temperature end of the onset curve was taken from the bound. Our
finite-temperature onset crosses $8/3$ only at $T/\sqrt B = 7.4\times10^{-4}$.

\subsection{The criterion is \texorpdfstring{$C_J/C_T < 8/(d^2(d+1))$}{C\_J/C\_T < 8/(d\^{}2(d+1))}}
\label{sec:4.4}

The coupling $\alpha = \kappa^2/\hat g^2$ depends on how the Yang--Mills term
is normalised relative to the Einstein--Hilbert term, and papers differ;
\cite{Wong:2013rda}, for instance, uses $2e^2L^2/\kappa^2$. A
convention-free form uses the central charges $C_T$ and $C_J$, which
normalise the two-point functions of the stress tensor and of the current
that the magnetic field couples to; criteria of this kind were written in
terms of $C_T$ and $C_J$ in~\cite{Nakayama:2015hga}, as an AdS form of the
weak gravity conjecture, for the $\mathrm{AdS}_2$ bound of the electric
brane. In the conventions
of~\cite{Osborn:1993cr}, Einstein gravity with a Maxwell term
$-F^2/4\hat g^2$ in $\mathrm{AdS}_{d+1}$ gives~\cite{Buchel:2009sk,Freedman:1998tz}
\begin{equation}
C_T = \frac{d+1}{d-1}\,\frac{\Gamma(d+1)}{\pi^{d/2}\,\Gamma(d/2)}\,\frac{L^{d-1}}{\kappa^2}, \qquad
C_J = \frac{(d-2)\,\Gamma(d)}{2\pi^{d/2}\,\Gamma(d/2)}\,\frac{L^{d-3}}{\hat g^2}~,
\end{equation}
where for the non-abelian field the second formula holds for each
generator, in particular for the $\tau^3$ that carries the background flux.
The ratio is
\begin{equation}
\frac{C_J}{C_T} = \frac{(d-1)(d-2)}{2d(d+1)}\,\alpha~,
\end{equation}
which is $3\alpha/20$ in $d = 4$, and the result of section~\ref{sec:4.3}
becomes
\begin{equation}
\text{the magnetised plasma is unstable at } T = 0 \quad\Longleftrightarrow\quad \frac{C_J}{C_T} < \frac{8}{d^2(d+1)}~.
\end{equation}
Here unstable means that the normal phase has a linear charged
instability at some field. $C_J$ is normalised by the algebra of the charges, $[Q^a, Q^b] = i\epsilon^{abc}Q^c$, so
that a doublet carries $T^3 = \pm\frac12$ (appendix~\ref{app:dictionary}).
The criterion holds for $d = 4, 5, 6$; in $d = 3$ the bound value $2/9$ is replaced by
$\alpha_c(3)/12 = 0.253$ (table~\ref{tab:dims}).

\begin{table}[t]
\centering
\small
\setlength{\tabcolsep}{4.5pt}
\begin{tabular}{cccccccc}
\toprule
$d$ & throat & $\alpha_\star(d)$ & $c_1/c_0$ & critical coupling & $C_J/C_T$ threshold & \multicolumn{2}{c}{supersymmetric check} \\
 & & & & & & $\alpha$ & $C_J/C_T$ \\
\midrule
$3$ & $\mathrm{AdS}_2 \times \mathbb{R}^2$ & $8/3$ & $+0.123$ & $\alpha_c(3) = 3.037$ & $0.253$ & $2$ & $1/6$ \\
$4$ & $\mathrm{AdS}_3 \times \mathbb{R}^2$ & $2/3$ & $-0.035$ & $\alpha_\star$ & $1/10$ & $1/4$ & $3/80$ \\
$5$ & $\mathrm{AdS}_4 \times \mathbb{R}^2$ & $4/15$ & $-0.350$ & $\alpha_\star$ & $4/75$ & -- & -- \\
$6$ & $\mathrm{AdS}_5 \times \mathbb{R}^2$ & $2/15$ & $-0.996$ & $\alpha_\star$ & $2/63$ & -- & -- \\
\bottomrule
\end{tabular}
\caption{The critical coupling by boundary dimension $d$: the throat, its
Breitenlohner--Freedman value $\alpha_\star(d) = 16/(d(d-1)(d-2))$, the
logarithmic ratio $c_1/c_0$ of the zero-temperature threshold mode,
$\rho^{(d-2)/2}w = c_0 + c_1\ln(\rho/\rho_\star)$ in the throat (positive signals a bound state in the
crossover), the resulting critical coupling, the threshold in $C_J/C_T$, and
the two supersymmetric theories against which the map is checked
(appendix~\ref{app:dictionary}).}
\label{tab:dims}
\end{table}

The criterion puts a number on the description
in~\cite{Ammon:2009xh,Wong:2013rda} of the coupling as the fraction of the
degrees of freedom that are charged. The one top-down case with a known
answer agrees with it. The $\mathcal{N} = 1$ R-current of any
four-dimensional superconformal theory has $C_R/C_T = 1/10$
(appendix~\ref{app:dictionary}), exactly the threshold for a vector of
unit R-charge. In $\mathcal{N} = 4$ super-Yang--Mills theory in an
equal-charge R-symmetry field, whose throat is that of D'Hoker and Kraus
uplifted on $S^5$, the $SO(6)$ $W$ bosons carry R-charge $4/3$ and so see
$\alpha = \tfrac23\big(\tfrac34\big)^2 = \tfrac38 < \alpha_\star$; their
throat mass $\ell_3^2m^2 = -4/3$, from eq.~(3.21) of~\cite{Donos:2011pn} at equal charges, is $-b$ at this
coupling, and they are the unstable mode found
in~\cite{Almuhairi:2011ws,Donos:2011pn}. That $3/8$ lies close to $\alpha_L$
does not place this theory in either lattice regime of section~\ref{sec:5},
which treats Einstein--$SU(2)$ theory only. In any theory whose magnetic
sector reduces to such a charged vector on an
$\mathrm{AdS}_{d-1} \times \mathbb{R}^2$ throat, $C_J/C_T$ below the
threshold implies an instability. The converse also needs the absence of a
bound state in the crossover, which we have shown for the
Einstein--Yang--Mills brane and which fails in $d = 3$.

We check the map from $\alpha$ to $C_J/C_T$ at two points where both sides
are fixed independently (appendix~\ref{app:dictionary}). In Romans' gauged
supergravity in $d = 4$, $\alpha_{\rm SUSY} = 1/4$, and the
$\mathcal{N} = 2$ Ward identities give $C_J(SU(2)_R)/C_T = 3/80 =
(3/20)(1/4)$. In minimal gauged supergravity in $d = 3$,
$\alpha_{\rm SUSY} = 2$ and $C_R/C_T = 1/6 =
\alpha/12$~\cite{Closset:2012ru}. Both points lie inside the
unstable window, the $d = 4$ one in the crystal regime of
section~\ref{sec:5}. Einstein--$SU(2)$ theory with $SU(2)$ flux is not a
consistent truncation of Romans' theory, since Romans' scalar is sourced
by the flux, so the comparison fixes the map, not the phase diagram of a string
vacuum. Whether any top-down theory reaches $\alpha_\star$ is open.

\subsection{No other charged mode is unstable above \texorpdfstring{$\alpha_\star$}{alpha*}}
\label{sec:4.5}

The critical coupling bounds the whole charged sector, not only the
polarised lowest Landau level for which it was derived. Static,
$z$-independent fluctuations of $W_\mu$ on any diagonal magnetic brane sort
by a single Landau index $n \ge 0$, since the Landau ladder operators shift
the level by one and the metric depends on $r$ only; $n = 0$ is the
polarised lowest Landau level (appendix~\ref{app:charged-channels}). Every
level with $n \ge 1$ is stable: its static energy, a sum of positive
radial terms and a positive semi-definite transverse form, vanishes only
on pure gauge, so no source-free, horizon-regular zero mode exists.
Only $n = 0$ has a negative transverse term, and its competition with the
radial gradient is the eigenvalue problem of section~\ref{sec:3}. The
argument uses only positivity, so it holds at every $\alpha$, $B$ and $T$,
and it survives momentum along $z$.

The argument also covers time-dependent perturbations. The brane has no
$A_t$ and no $g_{ti}$, so a growing mode must have real $\omega^2$, hence
$\omega = i\Gamma$ and negative static energy. An instability can therefore
only set in through a static zero mode at $\omega = 0$
(appendix~\ref{app:charged-channels}). Together with the $n = 0$
result, the charged sector has no linear instability above $\alpha_\star$.

We assume that the neutral sector, $A^3$ and the metric, is stable. No
growing neutral mode of the D'Hoker--Kraus brane is known. None has been
found at zero momentum in the scalar and tensor
channels~\cite{Janiszewski:2015ura,Shah:2024lym}, or, with a Chern--Simons
term, along the field~\cite{Ammon:2017ded}. Static modes across the field
respect the bound in the throat of the $U(1)^3$
truncation~\cite{Donos:2011pn}. In our theory, those of the sector the
condensate sources respect it (appendix~\ref{app:response-throat}) and have
no zero mode on any brane we solved (section~\ref{sec:5}). The other static
sectors, and all time-dependent modes across the field, have not been
computed. The known modulated instabilities of
magnetic branes need couplings our theory does not have: a neutral scalar
coupled to $F_{\mu\nu}F'^{\mu\nu}$, with $F'$ a second abelian field
strength~\cite{Donos:2011qt,Donos:2016hsd}, or, for the Chern--Simons
instability of~\cite{Nakamura:2009tf}, an electric background and a cubic
Casimir, which $SU(2)$ lacks.

\subsection{The approach to \texorpdfstring{$\alpha_\star$}{alpha*} has Berezinskii--Kosterlitz--Thouless form}
\label{sec:4.6}

We continued the onset calculation of section~\ref{sec:3} up the ladder of
section~\ref{sec:3.3} to $\beta = 10^{11}$. The interior reaches
the fixed point, with the invariant $\beta e^{-4V}$ at the horizon rising to
$6$, and the throat field $b_h$ at the horizon falls towards the bound
$1$ (appendix~\ref{app:numerics-onset}). Since
$2/(3b_h^2) = 6\alpha/(\beta e^{-4V})$ on the onset curve, together these
say $\alpha \to 2/3$. The approach is very slow: at $\beta = 10^{11}$,
where $B_c u_H^2 = 4.1 \times 10^5$, the coupling has only reached $0.591$.

This is the holographic Berezinskii--Kosterlitz--Thouless scaling
of~\cite{Kaplan:2009kr,Jensen:2010ga,Iqbal:2010eh,Jensen:2010vx,Evans:2010hi},
a Breitenlohner--Freedman bound crossed as a parameter is tuned. Here it
is the bound of an $\mathrm{AdS}_3$ throat, whereas
\cite{Iqbal:2010eh,Jensen:2010vx,Evans:2010hi} and the D3/D5 system
of~\cite{Jensen:2010ga} cross an $\mathrm{AdS}_2$ bound, and the holographic picture
of~\cite{Kaplan:2009kr} is a scalar below the bound of the full
$\mathrm{AdS}_{d+1}$. With $\nu = \sqrt{b_h - 1}$ the mode in the throat oscillates, up to
corrections of order $\varepsilon$, as
$\rho\,w \propto \cos(\nu\ln\rho + \ldots)$, with $\rho = r - r_0$, and the
onset is where the throat is long enough to hold the solution that joins
the horizon to the ultraviolet without an extra node.

Both ends of the oscillation can be computed, and matching them has no free
parameter. At the horizon the solution on the BTZ throat is a
hypergeometric function, whose connection formula gives a phase $\chi(\nu)$
in closed form. At the ultraviolet end, the outer region of every rung is,
in scaled variables, the same zero-temperature brane; integrating the
source-free mode inwards from the boundary on it gives a phase
$\phi_\star(\nu)$ in the throat. The two solutions match when
\begin{equation}
\nu\,\ln(\rho_\star/h) = \chi(\nu) + \phi_\star(\nu)~,
\end{equation}
where $h$ is the BTZ horizon and $\rho_\star$ the radius at which the
proper area $e^{2V}$ of the flux plane has doubled from its throat value
(figure~\ref{fig:matching}). This condition reproduces $\nu$ on the ladder
and the direct solves by horizon shooting (appendix~\ref{app:numerics-matching}).

\begin{figure}[t]
\centering
\includegraphics{matching_diagram}
\caption{The matching diagram: the marginal mode on the $\beta = 10^8$ rung
against $\ln(\rho/h)$, with the exact BTZ horizon solution and the $T = 0$
outer mode laid over it, the three regions shaded and $h$, $\rho_\star$ and
the crossover radius $r_c$ marked.}
\label{fig:matching}
\end{figure}

As $\nu \to 0$, $\chi + \phi_\star \to \pi$: the throat holds one
half-oscillation of the mode. With $\ln(\rho_\star/h) \simeq \tfrac14\ln\beta$
(appendix~\ref{app:numerics-matching}),
$\nu^2 \simeq (\alpha_\star - \alpha)/2\alpha_\star$ and
$\beta \propto B^2/T^4$ at fixed field, the onset obeys
\begin{equation}
\ln\beta_c \simeq \frac{8\pi}{\sqrt{3(\alpha_\star - \alpha)}}~, \qquad
T_c \propto \sqrt{B}\,\exp\!\Big(-\frac{2\pi}{\sqrt{3(\alpha_\star - \alpha)}}\Big) = \sqrt B\,e^{-3.63/\sqrt{\alpha_\star-\alpha}}~.
\end{equation}
This is the law of~\cite{Iqbal:2010eh,Jensen:2010vx}; the same law arises
from the annihilation of fixed points, with no bound crossed,
in~\cite{Faedo:2019nxw}. Since $\alpha > \alpha_L$ wherever it applies,
$T_c$ here is the spinodal of the normal phase (section~\ref{sec:5.3}): the
essential singularity is that of the onset of the instability, not of an
infinite-order transition.

The law is asymptotic, and it sets in only at extremely low temperature.
The horizon phase is analytic in $\nu$,
$\chi = \pi/2 - 2\nu\ln 2 + O(\nu^3)$, so its corrections only shift
$\ln\beta_c$ by a constant. The outer phase is not. At $\alpha_\star$ the
mode arrives in the throat almost entirely as the non-logarithmic solution,
with $c_1/c_0 = -0.035$, and matching it onto
$\cos(\nu\ln(\rho/\rho_\star) + \phi_\star)$ gives
$\tan\phi_\star = |c_1/c_0|/\nu + 0.50\,\nu + O(\nu^3)$. So $\phi_\star$
reaches its limit $\pi/2$ only for $\nu \ll 0.035$, that is
$\ln\beta \gg 400$. Over the throat part of the ladder, $\beta \ge 10^3$ or
$0.25 \le \nu \le 0.6$, $\phi_\star$ instead sits on a plateau of about
$0.22$--$0.25$ (figure~\ref{fig:approach}, right). In this range the
throat needs to hold less phase, so the plasma condenses more easily than the asymptotic law
predicts. A direct fit on the ladder gives
\begin{equation}
T_c \propto \sqrt B\,e^{-2.25/\sqrt{\alpha_\star - \alpha}}~,
\end{equation}
a coefficient $38\%$ below the asymptotic one. Shooting from the horizon
far beyond the ladder confirms that the effective coefficient rises
slowly towards its asymptotic value, as the matching condition predicts
(appendix~\ref{app:numerics-matching}); the asymptotic law takes over only
at $T_c/\sqrt B \sim e^{-100}$.

\begin{figure}[t]
\centering
\includegraphics{approach_law}
\caption{The approach law. Left: $1/\nu$ against $\ln\beta$ with the
ladder, the matching prediction, the ladder fit and (inset, to
$\ln\beta = 2000$) the $1/4\pi$ asymptote. Right: $\chi$, $\phi_\star$ and
$\chi + \phi_\star$ against $\nu$ with the $\pi$ limit, the turnover at
$\nu = |c_1/c_0| = 0.035$ and the throat part of the ladder, $0.25 \le \nu \le 0.6$.
Below $\nu = 0.02$ (dashed) $\phi_\star$ is the small-$\nu$ form
$\pi/2 - \arctan(\nu|c_0/c_1|)$ rather than a direct solve, and so is the
approach to $1/4\pi$ in the inset.}
\label{fig:approach}
\end{figure}
